\section{Terms That Must Not Be Collapsed}

\subsection{Chastity: the governing virtue}

Chastity is not a synonym for abstinence and not a peculiarly clerical state. It is the moral virtue that integrates sexuality within the truth of the person, covenant, and vocation. The married person lives chastity through faithful conjugal love, including morally ordered marital intercourse. The unmarried person lives chastity through sexual abstinence. A violation of chastity is therefore possible in every state, but the concrete acts required by chastity differ by state.

Clerical status adds obligations and representative significance; it does not erase the natural and sacramental reality of a valid marriage. CCEO canon 374 consequently speaks of both celibate clerics and married clerics shining in the splendor of chastity. The 1998 directory for permanent deacons similarly treats a married deacon's sacramental marriage and family life as a genuine field of sanctification rather than an embarrassment tolerated alongside ministry.

\subsection{Continence: abstention from genital acts}

Continence denotes refraining from genital acts. It may be temporary, periodic, or perfect and perpetual. First Corinthians 7:5 describes temporary abstinence by mutual consent for prayer; that is not celibacy. CIC canon 277 \S1 joins ``perfect and perpetual continence for the sake of the kingdom of heaven'' to the clerical obligation of celibacy. The precise application of that canon to married permanent deacons is examined separately, because the code's general wording and the Church's authoritative postconciliar implementation have generated a genuine interpretive dispute.

The phrase \emph{perfect continence} in clerical law means complete sexual abstinence, not merely fidelity to one spouse. The adjective cannot be weakened to mean ordinary marital chastity. The legal question is instead \emph{who is bound by that form of continence, by what norm, and from what moment?}

\subsection{Celibacy: an unmarried juridic condition}

Celibacy is the condition of not being married. In Latin clerical formation it is also a freely assumed, public ecclesial commitment made before diaconal ordination by a candidate who is not already bound by a public perpetual vow of chastity in a religious institute (CIC can. 1037). It is not normally a religious vow for a diocesan cleric. A religious cleric may additionally be bound by the public vow of chastity proper to consecrated life (CIC can. 599); the vow and the clerical obligation have different juridic titles even where their practical sexual demands coincide.

A widower is unmarried and thus celibate in an ordinary descriptive sense, but his canonical history differs from that of a man never married. A married permanent deacon is a cleric without being celibate. A laicized priest may remain bound by celibacy unless the Roman Pontiff grants the separate dispensation contemplated by CIC canon 291. Precise language prevents pastoral harm.

\subsection{Four further distinctions}

\begin{fourstable}{Distinction}{Do not collapse}{Legal consequence}{Pastoral consequence}
Ordination and marriage & Priesthood does not by sacramental nature require an unmarried recipient. & Vatican II, \emph{Presbyterorum ordinis} 16, expressly says celibacy is not demanded by the nature of priesthood. & Married Eastern priests and married former non-Catholic ministers are not lesser priests.\\
Selection and validity & A discipline governing who may be ordained is not necessarily a condition for valid ordination. & Breach may concern liceity, impediment, irregularity, or governance without making the sacrament invalid. & Concrete cases require the exact canon and facts, not slogans.\\
Marriage before and after orders & Admission of a married man is not permission for a cleric to contract marriage. & Sacred orders is a diriment impediment in both codes (CIC can. 1087; CCEO can. 804). & A wife dies, civil divorce occurs, or a declaration of nullity is issued: none by itself grants permission to marry.\\
Discipline and doctrine & A mutable ecclesiastical discipline is not meaningless or devoid of doctrinal roots. & The Latin norm can change; the Church authoritatively teaches its scriptural, Christological, ecclesial, and eschatological fittingness. & One may discuss reform without denying either the legitimacy of married clergy or the value of celibacy.\\
\end{fourstable}

\Needspace{12\baselineskip}
\subsection{A compact status map}

\begin{threestable}{Clerical state}{Ordinary sexual discipline}{Juridic key}
Latin bishop or presbyter selected under the ordinary norm & Celibate chastity and perfect, perpetual continence. & CIC cann. 277, 1037; episcopal selection presupposes priesthood and the Latin norm.\\
Married Latin permanent deacon & Conjugal chastity within a valid marriage; no ordinary right to marry after ordination. & CIC cann. 1031 \S2, 1050, 1087; authoritative diaconal norms and the can. 277 question discussed below.\\
Unmarried Latin permanent deacon & Public assumption of celibacy before ordination; perfect, perpetual continence. & CIC can. 1037.\\
Eastern Catholic bishop & Candidate must not be bound by marriage. & CCEO can. 180, 3\textsuperscript{o}.\\
Eastern Catholic married presbyter or deacon & Conjugal chastity; admission and formation governed by common and applicable particular law. & CCEO cann. 373--375, 758 \S3, 769 \S1, 2\textsuperscript{o}.\\
Eastern Catholic celibate presbyter or deacon & Celibate chastity according to common and particular law. & CCEO cann. 373--374.\\
Any cleric seeking to marry after ordination & Marriage is invalidly attempted absent a dispensation from the impediment by competent authority where the law permits it. & CIC can. 1087; CCEO can. 804; penal and status consequences may also follow.\\
\end{threestable}
